% Part: applied-modal-logics
% Chapter: epistemic-logic
% Section: introduction

\documentclass[../../../include/open-logic-section]{subfiles}

\begin{document}

\olfileid[tr]{aml}{el}{int}

\olsection{Giriş}

Modal mantık nasıl \emph{modal önermeleri} ve bunlar arasındaki gerektirme
bağıntılarını ele alıyorsa epistemik mantık da \emph{epistemik önermeleri} ve
bunlar arasındaki gerektirme bağıntılarını ele alır. Birli bağlaçlar, modal
işleçleri olanak ve zorunluluğun temsilleri olarak yorumlamak yerine bilgi ve
inancı modelleyecek biçimde epistemik ya da doksastik olarak yorumlanır.
Örneğin aşağıdaki gibi iddiaları ifade etmek isteyebiliriz:

\begin{enumerate}
\item Richard, Calgary'nin Alberta'da olduğunu bilir.
\item Audrey, kanepede bir köpek bulunmasının mümkün olduğunu düşünür.
\item Richard, Audrey'nin dersinin salı günleri olduğunu bildiğini bilir.
\item Herkes bir yılda 12 ay olduğunu bilir.
\end{enumerate}
Çağdaş epistemik mantığın başlangıcı çoğu kez Jaakko Hintikka'nın 1962 tarihli
\emph{Knowledge and Belief} adlı eserine götürülür; eser, olası dünyalar
semantiğinin mantıkta giderek daha yaygın kullanıldığı bir dönemde yazılmıştır.
% [tr source emendation] Upstream misspells Hintikka's first name as `Jaako`;
% the established spelling is `Jaakko`.
Epistemik mantıklar gerçekte diğer modal mantıklarla aynı semantik araçların
çoğunu kullanır, fakat onları farklı yorumlar. Başlıca değişiklik,
\emph{erişilebilirlik bağıntısının} neyi temsil ettiğine ilişkindir. Epistemik
mantıklarda bu bağıntılar bir tür \emph{epistemik olanaklılığı} temsil eder.
Modellediğimiz epistemik kavramın, erişilebilirlik bağıntısına koymak
istediğimiz kısıtları etkilediğini göreceğiz. Karşılıklılık kuramına epistemik
bir yorum verildiğinde ne olduğunu da inceleyeceğiz. Yukarıdaki örneklerin iki
özneden, Richard ile Audrey'den, ve her birinin bildikleri arasındaki ilişkiden
söz ettiğine dikkat edin. Ele alacağımız epistemik mantıklar, bu tür şeylerin
ifade edilebildiği çok özneli mantıklar olacaktır. Buna karşılık tek özneli bir
epistemik mantık yalnızca bir bireyin ne bildiğinden ya da neye inandığından
söz eder.

\end{document}
